\chapter{私募股权与另类投资版图}

\section{另类投资在资产配置中的位置}

另类投资（Alternative Investments）泛指在传统股票与债券之外、流动性通常较低、结构更复杂、信息披露与监管框架与公募不同的资产类别。机构组合引入另类资产的核心动机包括：\textbf{分散化}（与公开市场相关性不完全同步）、\textbf{风险溢价}（流动性折价与复杂度溢价）、以及在某些市场环境下\textbf{绝对收益}或\textbf{通胀对冲}的诉求。私募股权（PE）是其中体量最大、制度最成熟的一条主线，但与对冲基金、实物资产、私募信贷等边界常有交叉。

\section{私募股权的主要形态}

\textbf{并购基金（Buyout）}通常以控股或重大少数股权方式收购成熟企业，依托杠杆、运营改善与行业整合提升退出倍数。\textbf{成长资本（Growth）}面向已盈利或接近盈利的扩张期企业，股权稀释程度与估值弹性介于 VC 与并购之间。\textbf{风险投资（VC）}聚焦早期技术与商业模式验证，单笔损失率可较高，依赖幂律分布中的极端成功项目。\textbf{夹层资本（Mezzanine）}与\textbf{困境/特殊机会}策略则更接近信贷或混合工具，在资本结构中的位置与纯股权 PE 不同。读者在本书语境下可将「PE」理解为\textbf{以非上市股权或相关工具为主、以主动治理与退出为收益来源}的广义家族，而不局限于某一细分标签。

\section{与公募、对冲基金的对照}

\textbf{期限与流动性：}典型 PE 基金有固定存续期与资本召唤（capital call）节奏，LP 在基金层面难以按净值每日申赎；公募与多数对冲基金则强调份额流动性（尽管侧袋、门控等机制会削弱这一特征）。\textbf{定价：}非上市股权缺乏连续报价，估值依赖季度评估与交易事件，业绩归因与「账面回报」解读需要额外谨慎。\textbf{监管：}面向零售的公募受最严格的信息披露与销售规则约束；私募链条在各国法域下对应不同「合格投资者」门槛与募集方式限制（详见本书合规篇）。\textbf{收益结构：}PE 更强调\textbf{控制或强影响}下的价值创造叙事；多策略对冲基金可能以市场中性、套利或宏观为主，与单一企业基本面绑定度较低。

\section{为何「版图」思维重要}

从零到一设立或参与一支基金时，先明确自己在版图中的\textbf{坐标}（阶段、行业、地域、币种、是否杠杆），能避免基金条款与团队能力与宣称策略错位。后续章节中的募资对象、尽调深度、条款谈判与退出假设，都隐含了这一坐标选择；合规篇则进一步约束该坐标在特定司法辖区是否可合法落地。
